\section{Identifiability Diagnostics: Definitions}
\label{app:diagnostics_definitions}

This appendix gives the full definitions of the diagnostics summarized in
Section~\ref{sec:identifiability_theory} and Table~\ref{tab:diagnostics}.
It uses the general notation of Section~\ref{sec:problem_setup}; ``background'' refers to
the nuisance basis \(Q\), and a ``coefficient'' is one component \(j\).

\begin{table}[t]
\centering
\small
\setlength{\tabcolsep}{3pt}
\caption{Identifiability diagnostics and reporting outputs produced with source
activity estimates.  The lower block lists the set- and graph-valued outputs that
accompany the scalar diagnostics.}
\label{tab:diagnostics}
\begin{tabularx}{\columnwidth}{l
  >{\raggedright\arraybackslash}X
  >{\raggedright\arraybackslash}X
  >{\raggedright\arraybackslash}X}
\toprule
Symbol & Name & Interpretation & Action \\
\midrule
\(r_{\mathrm{num}}(\widetilde A)\) & numerical rank & exact separability & require rank \(J\) \\
\(\sigma_J(\widetilde A)\) & padded min.\ singular value & noise robustness & flag zero/small \\
\(\kappa(\widetilde A)\) & condition number & error amplification & \(\infty\) if deficient \\
\(r_{\mathrm{eff}}\) & effective rank & noise-level resolution & reduce grouping \\
\(v_j\) & coefficient visibility & measurable effect & flag weak \\
\(a_j\) & background absorption & signal lost to \(Q\) & warn after projection \\
\(\rho_{ij}\) & pairwise coherence & pairwise similarity & merge if \(>\tau_\rho\) \\
\(d^{\mathrm{ray}}_{ij}\) & ray distance & \(\sqrt{1-\rho_{ij}^2}\) & N/A if weak \\
\(s_g\) & separation & noise-level separability of \(g\) & report \(g\) if \(\ge\tau_\theta\) \\
\midrule
\(\mathcal W\) & weak set & undetectable coefficients & flag \\
\(\mathcal P^\star\) & matroid components & exact resolution & atoms with \(\mathcal D\) \\
\bottomrule
\end{tabularx}
\end{table}

For floating-point computation, the default numerical rank tolerance is
\(\tau_{\mathrm{num}}=\max(N,J)\,\epsilon_{\mathrm{mach}}\,\sigma_1\) with
numerical rank \(r_{\mathrm{num}}=\#\{i:\sigma_i>\tau_{\mathrm{num}}\}\),
where \(\epsilon_{\mathrm{mach}}\) is the machine precision of the SVD dtype
(about \(2.2\times10^{-16}\) in double precision).  This
\(\tau_{\mathrm{num}}\) is a purely numerical tolerance for detecting exact rank
deficiency; it is distinct from the scientific, noise-dependent threshold
\(\tau_\sigma\) below.  The primary identifiability score is the
\emph{padded minimum singular value} \(\sigma_J(\widetilde A)\), where the
singular spectrum is padded with zeros whenever \(J>N\) or
\(\widetilde A\) is rank deficient.  A rank-deficient response therefore
scores zero and is never reported as stable.
The condition number is \(\kappa(\widetilde A)=\sigma_1/\sigma_J\) when
\(r_{\mathrm{num}}=J\) and \(\infty\) when \(r_{\mathrm{num}}<J\).
The smallest numerically positive value may additionally be reported as
\(\sigma_{\min,+}\), but it is not the identifiability score and never turns a
rank-deficient matrix into a stable one.  For a separately chosen,
noise-dependent threshold \(\tau_\sigma\), the effective rank is
\(r_{\mathrm{eff}}(\tau_\sigma)
=\#\{i:\sigma_i(\widetilde A)>\tau_\sigma\}\).
Unlike \(\tau_{\mathrm{num}}\), the threshold \(\tau_\sigma\) is meant to be set
from the observation-noise scale, not from machine precision; it is declared before
fitting and is not tuned on source-recovery results.  When no \(\tau_\sigma\) is
declared, the implementation sets \(\tau_\sigma=\tau_{\mathrm{num}}\), so that
\(r_{\mathrm{eff}}\) equals the numerical rank; this is the case in every reported
run.

For coefficient-specific diagnostics, let \(\widetilde{\mathbf a}_j\),
\(j=(k,b)\), be a projected source--basis fingerprint.  Its visibility and the
weak set at threshold \(\tau_v\ge0\) are
\(v_j=\|\widetilde{\mathbf a}_j\|_2\) and
\(\mathcal W=\{j:v_j\le\tau_v\}\).
A weak coefficient has little measurable effect after transport and background
correction: \(v_j\) is the sensor-time magnitude produced by one unit of
coefficient \(c_j\), and \(\tau_v\) marks fingerprints below the minimum
detectable signal.  A practical predeclared rule sets \(\tau_v\) from the noise
scale and a target signal-to-noise ratio, e.g.\ \(\tau_v=\mathrm{SNR}_{\min}\,
\sigma_e\), so a unit coefficient is flagged weak when even its full response is
indistinguishable from noise.  Every reported run uses \(\tau_v=0\), so only an
exactly zero fingerprint is weak.  For eligible indices \(i,j\notin\mathcal W\),
pairwise ambiguity is measured by projected fingerprint coherence
\(\rho_{ij}=|\widetilde{\mathbf a}_i^\top\widetilde{\mathbf a}_j|/
(\|\widetilde{\mathbf a}_i\|_2\|\widetilde{\mathbf a}_j\|_2)\).
Values near one indicate nearly proportional sensor-time signatures.  The
ambiguous-pair set is
\(\mathcal A=\{(i,j):i,j\notin\mathcal W,\ \rho_{ij}>\tau_\rho\}\),
\(\tau_\rho\in[0,1]\).  The threshold \(\tau_\rho\) is predeclared, not optimized after inspecting
source recovery.  Coherence and ray distance are
reported only for eligible (nonweak) pairs; entries involving a weak fingerprint
are shown as ``N/A'' and excluded from extrema, while the weak-visibility flag
is always reported.

\noindent\textbf{Near-dependencies and \(\sigma_J\).}  If a unit vector \(\mathbf w\)
gives \(\|\sum_j w_j\widetilde{\mathbf a}_j/\|\widetilde{\mathbf a}_j\|_2\|_2=\delta\), then
\(\mathbf v=(w_j/\|\widetilde{\mathbf a}_j\|_2)_j\) has \(\|\widetilde A\mathbf v\|_2=\delta\) and
\(\|\mathbf v\|_2\ge1/\max_j\|\widetilde{\mathbf a}_j\|_2\), so
\(\sigma_J\le\delta\max_j\|\widetilde{\mathbf a}_j\|_2\).  A near-dependency among three or
more columns with no high pairwise coherence therefore lowers \(\sigma_J\) without
triggering the pairwise merge rule.

\noindent\textbf{Classical connections.}  With
\(\widetilde{\boldsymbol\epsilon}=P_Q^\perp\boldsymbol\epsilon\) and
\(\boldsymbol\epsilon\sim N(0,\sigma^2I)\), the Fisher information for \(\mathbf c\) is
\(\widetilde A^\top\widetilde A/\sigma^2\), because
\(\operatorname{col}(\widetilde A)\perp\operatorname{col}(Q)\), so the rank condition of
Proposition~\ref{prop:rank_identifiability} is a nonsingular Fisher information
\citep{rothenberg1971identification}, and a linear combination
\(\boldsymbol\lambda^\top\mathbf c\) is estimable exactly when
\(\boldsymbol\lambda\in\operatorname{row}(\widetilde A)\) \citep{searle1971linear}.  When
only components \(i\) and \(j\) are fitted, the correlation between their least-squares
estimates is \(-\widetilde{\mathbf a}_i^\top\widetilde{\mathbf a}_j/
(\|\widetilde{\mathbf a}_i\|_2\|\widetilde{\mathbf a}_j\|_2)\), so \(\rho_{ij}\) is its
magnitude.

\noindent\textbf{Visibility under operator error.}  If
\(\widetilde A'=\widetilde A+\Delta\), the triangle inequality gives
\(|\,\|\widetilde{\mathbf a}'_j\|_2-\|\widetilde{\mathbf a}_j\|_2|\le\|\Delta\mathbf e_j\|_2\),
so the weak set is unchanged when \(|v_j-\tau_v|>\|\Delta\mathbf e_j\|_2\) for every \(j\).

To express the same scale-invariant geometry as a distance, we also report the
ray distance for eligible, generally signed projected fingerprints,
\(d^{\mathrm{ray}}_{ij}
=\min_{\alpha\in\mathbb R}
\|\widetilde{\mathbf a}_i-\alpha\widetilde{\mathbf a}_j\|_2/
\|\widetilde{\mathbf a}_i\|_2\).
For nonweak fingerprints,
\begin{equation}
d^{\mathrm{ray}}_{ij}=\sqrt{1-\rho_{ij}^2},
\label{eq:ray_distance_coherence_relation}
\end{equation}
which we derive in Appendix~\ref{app:identifiability_proofs}; numerical
implementations evaluate \(\sqrt{\max(0,1-\rho_{ij}^2)}\).  This is an unoriented
linear-ray diagnostic, not a claim that projected fingerprints are nonnegative.
High coherence and small ray distance encode the same pairwise geometry, so ray
distance is a diagnostic geometrically equivalent to coherence under this
normalization, not an independent identifiability criterion.  Like coherence, it
is ``N/A'' for pairs involving \(\mathcal W\) and excluded from distance extrema.

\noindent\textbf{Background absorption.}
Background correction can improve robustness by removing broad temporal
variation, but it can also remove source-driven signal.  Let \(P_Q=QQ^\dagger\).
For coefficient fingerprint \(j=(k,b)\), define the background absorption
fraction \(a_j=\|P_Q\mathbf a_j\|_2/\|\mathbf a_j\|_2\).
If \(\|\mathbf a_j\|_2\) is zero at numerical tolerance,
\(a_j\) is undefined and reported as N/A; the coefficient remains explicitly
weak.
If \(a_j\approx 1\), most of the raw source--basis fingerprint lies in the
background space, and the projected visibility \(v_j\) will be small even if
the raw fingerprint is large.  This is why identifiability must be assessed
after projection.


\section{Uncertainty, Adequacy, and Reporting Details}
\label{app:reporting_details}

This appendix expands the reporting components summarized in
Section~\ref{sec:identifiability_theory}, and gives the IASA thresholds
(Table~\ref{tab:thresholds}).

\begin{table}[t]
\centering
\small
\setlength{\tabcolsep}{3pt}
\caption{IASA thresholds, their meaning, and how they are set.  None is
optimized on source-recovery error.}
\label{tab:thresholds}
\begin{tabular}{lp{0.3\linewidth}p{0.42\linewidth}}
\toprule
Threshold & Meaning & How set \\
\midrule
\(\tau_{\mathrm{num}}\) & numerical rank tol.\ & \(\max(N,J)\epsilon_{\mathrm{mach}}\sigma_1\) \\
\(\tau_\sigma\) & effective-rank floor & declared noise scale; \(\tau_{\mathrm{num}}\) if none \\
\(\tau_v\) & visibility floor & e.g.\ \(\mathrm{SNR}_{\min}\,\sigma_e\); \(0\) in every reported run \\
\(\tau_\theta\) & separation threshold & \(\sqrt{1-0.99^2}=0.141\) in every reported run, or \(\sigma(\sqrt{\operatorname{rank}\widetilde A}+\sqrt{2\log(1/\delta)})/\eta\) \\
\(\tau_\rho\) & pairwise coherence & \(0.99\); for two columns \(\rho>\tau_\rho\) exactly when \(s<\tau_\theta\) \\
flag & identifiability flag & merge, nonempty \(\mathcal W\), visibility ratio \(\le0.1\), absorption \(\ge0.9\), \(\sigma_J\le10^{-6}\), or coherence \(\ge0.99\) \\
\(\epsilon_w\) & wind-drift cap & physical wind tolerance \\
\bottomrule
\end{tabular}
\end{table}

\noindent\textbf{Uncertainty reporting.}
Let \(\widetilde r=\widetilde Y-\widetilde A\widehat{\mathbf c}\) be the
projected residual and
\(\widehat\sigma^2=\|\widetilde r\|_2^2/\max(N-r_{\mathrm{eff}}(\tau_\sigma),1)\)
a residual variance estimate.  For an unconstrained or active-set approximation,
let \(\mathcal A_c=\{j:\widehat c_j>\tau_c\}\) and let
\(\widetilde A_{\mathcal A_c}\) contain the active columns.  The covariance
of the active estimates is approximated by
\begin{equation}
\operatorname{Cov}(\widehat{\mathbf c}_{\mathcal A_c})
=
\widehat\sigma^2
(\widetilde A_{\mathcal A_c}^\top
\widetilde A_{\mathcal A_c})^\dagger.
\label{eq:active_covariance}
\end{equation}
The use of \(N-r_{\mathrm{eff}}(\tau_\sigma)\) is a heuristic degrees-of-freedom
count that does not exactly account for the background rank \(r\) or the NNLS
active-set constraints.  Equation~\ref{eq:active_covariance} is the standard
active-set least-squares covariance approximation
\citep{seber2003linear}: its diagonal entries give per-coefficient uncertainty
and its off-diagonal entries expose coupled (jointly uncertain) estimates.  For
each transport ensemble member from
Appendix~\ref{subsec:wind_field_estimation}, IASA rebuilds
\(\widetilde A^{(b)}\), refits \(\widehat{\mathbf c}^{(b)}\), and reports
empirical quantiles across members as transport uncertainty.  Every ensemble is
tagged by provenance: wind and dispersion members have
\texttt{ensemble\_kind=transport} and may form transport-uncertainty intervals,
whereas alternative inventory maps, source locations, labels, or scales have
\texttt{ensemble\_kind=inventory} and produce scenario-wise robustness tables.
The two kinds are never pooled into one interval, and active-set intervals
describe the fitted observation model without converting inventory scenarios into
probabilistic uncertainty.

\noindent\textbf{Residual model-adequacy check.}
Let \(Z\) be an orthonormal basis for the background-orthogonal observed-row
space, and define the nondegenerate residual coordinates and externally supplied
noise covariance
\begin{equation}
\bar r=Z^\top(Y-A\widehat{\mathbf c}),
\qquad
\bar\Sigma_e=Z^\top\Sigma_e Z,
\label{eq:adequacy_coordinates}
\end{equation}
where \(\Sigma_e\) is an externally declared observation-noise covariance and
\(Z\) removes the background directions so that \(\bar r\) carries no degrees of
freedom absorbed by \(Q\).  For independent homoscedastic noise
\(\Sigma_e=\sigma_e^2 I\), orthonormality of \(Z\) gives
\(\bar\Sigma_e=\sigma_e^2 I\) and \(T_{\mathrm{res}}=\|\bar r\|_2^2/\sigma_e^2\).
When \(\bar\Sigma_e\) is calibrated independently of this fitted residual and is
positive definite, the adequacy statistic is
\begin{equation}
T_{\mathrm{res}}=\bar r^\top\bar\Sigma_e^{-1}\bar r.
\label{eq:residual_adequacy_statistic}
\end{equation}
Its null distribution is not a plain \(\chi^2\), because \(\widehat{\mathbf c}\)
is fitted from the same data; IASA therefore calibrates it by parametric
bootstrap.  It simulates \(B_{\mathrm{res}}=1000\) synthetic observation sets from
the fitted source/background model plus the declared noise model, refits the
coefficients on each draw, and recomputes \(T_{\mathrm{res}}\) to form the null
distribution.  At \(\alpha=0.05\) it flags inadequacy when the observed statistic
exceeds the empirical \((1-\alpha)\)-quantile; the Monte Carlo \(p\)-value uses
the standard add-one correction.  If no externally calibrated noise model is
supplied, IASA reports raw and projected residual norms and sensor-wise,
time-wise, and autocorrelation summaries with
\texttt{calibration\_status=uncalibrated} and emits no pass/fail adequacy claim.
Rejection establishes a mismatch with the declared model but does not identify a
missing source, and non-rejection does not establish inventory completeness: an
omitted source whose sensor-time signature lies in
\(\operatorname{span}[A,Q]\) can be absorbed by fitted terms
and remain residual-invisible.

\noindent\textbf{Wind-distribution diagnostics.}
Diagnostics for one realized wind window answer a retrospective question.
Prospective network adequacy is assessed empirically over predeclared historical
or simulated wind windows.  Across those response matrices, IASA reports the
5th, 50th, and 95th percentiles of \(\sigma_J\), the probability of full
numerical and effective rank, the probability that each coefficient is weak,
and the frequency of each grouping
\(\mathcal P^\star\vee\mathcal D\).  These are Monte Carlo or empirical
distributional summaries under the sampled wind population, not a new
identifiability theorem.

\noindent\textbf{Acceptance criteria for refinement.}
If the optional constrained refinement of
Appendix~\ref{app:transport_response} is used, it is accepted only if it improves
fit without degrading separability.  Let \(\widetilde A_{0}\) be the
projected response before refinement and \(\widetilde A_{\mathrm{ref}}\)
after refinement.  We require
\(\sigma_J(\widetilde A_{\mathrm{ref}})\ge
(1-\eta_{\mathrm{id}})\sigma_J(\widetilde A_{0})\) and
that the grouping \(\mathcal P^\star\vee\mathcal D\) is unchanged.  These checks prevent end-to-end tuning from improving
sensor fit by making source--basis fingerprints less distinguishable.  If either
response has fewer than \(J\) numerically nonzero singular values, its
\(\sigma_J\) is taken to be zero
(Section~\ref{sec:identifiability_theory}), so refinement cannot be
accepted by exploiting a rank-deficient response.

\noindent\textbf{Reporting protocol.}
The atoms are the blocks of \(\mathcal P^\star\vee\mathcal D\), with \(\mathcal P^\star\)
computed at \(\tau_{\mathrm{num}}\) and \(\mathcal D\) the declared source partition.  The
reported partition is the partition into unions of atoms with every block's separation
at least \(\tau_\theta\) that has the most blocks, with ties broken by the largest
smallest separation and remaining ties listed (Section~\ref{sec:method}).  For each block
\(G\), IASA reports the group signal \(\At_G\widehat{\mathbf c}_G\), its \(L_1\) share, its
separation \(s_G\), the bound \(b_G\) of Theorem~\ref{thm:group_error}(c) at the declared
noise level, the member visibilities and, for a source with several basis columns, the
activity trajectory \(\widehat{\boldsymbol\theta}_{G,1:T}\).  Coefficients are reported only
for blocks with one component.  Weak coefficients keep their basis names in the report,
and coherences involving weak fingerprints are displayed as ``N/A''.  The recorded runs
computed their groupings with the pairwise coherence rule at \(0.99\) (merge two
declared blocks when a pair of their nonweak columns exceeds it); on every saved
operator and observed week this gives the same partition as the separation rule.
